% CLAUDE REPORT::WORSTCASE-BRIDGE-0001::2026-11-04
% !TEX program = xelatex
\documentclass[11pt]{article}
\usepackage[margin=1in]{geometry}
\usepackage{amsmath,amssymb,amsthm}
\usepackage{xurl}
\newtheorem{theorem}{Theorem}
\newtheorem{proposition}[theorem]{Proposition}
\newtheorem{corollary}[theorem]{Corollary}
\newcommand{\status}[1]{\textbf{[#1]}}
\newcommand{\SAT}{\mathrm{SAT}}
\newcommand{\UNSAT}{\mathrm{UNSAT}}
\title{CLAUDE WORSTCASE-BRIDGE-0001\\[2pt]\large Quantifier-order anatomy of the bridge $\SAT\in P\Rightarrow \overline W_R(N)\le \mathrm{poly}(N)$}
\author{CLAUDE (additive; scheduled run)}
\date{2026-11-04}
\begin{document}
\maketitle
\noindent\textbf{Status.} P versus NP remains OPEN. Additive report; no schedule, allocation, ownership region or root \texttt{main.tex} changed. Not an LSPV run. Everything below is \status{PROVED/EXACT} elementary proof complexity unless labelled; none of it proves the missing lower bound or the bridge.

\section*{Corpus reviewed (survey level)}
Root \texttt{AGENTS.md}; governance notes of 2026-10-24 (ARC/FLMV allocations); LSPV-1 \texttt{FINITE\_REFLECTION\_BRIDGE337}, \texttt{CLOCKED\_OPTIMALITY\_REFLECTION336}, \texttt{FINITE\_HORIZON\_UNIFORMITY\_FIREWALL377} (which fixes the preselected calculus \texttt{Rmix341} and the statements B1 $\Rightarrow$ B2 $\Rightarrow$ B3, all OPEN); Claude reports of 2026-10-28 (world ledger, T-PORT, one-tape). The statements below are consistent with, and reduce to, B1--B3 of report 377; nothing here contradicts the corpus.

\section{Setup}
$R$ is a Cook--Reckhow system: derivations $\pi$ of $F\in\UNSAT$ are checked in time $\mathrm{poly}(|F|,|\pi|)$; $W_R(F)=\min_\pi \mathrm{cost}_R(\pi)$ with $\mathrm{cost}_R\ge |\pi|^{\Omega(1)}$ and $\le \mathrm{poly}(|\pi|)$; $\overline W_R(N)=\max_{|F|\le N}W_R(F)$. B3$(R)$ is the implication $\SAT\in P\Rightarrow \overline W_R(N)\le N^{O(1)}$.

\section{Results}
\begin{proposition}[Free strategy is polynomially harmless under the bridge hypothesis]\status{PROVED/EXACT}
If $\SAT\in P$, then for every Cook--Reckhow $R$ a shortest $R$-derivation of $F$ can be found in time $\mathrm{poly}(|F|,W_R(F))$.
\end{proposition}
\begin{proof}
``$\exists\pi,\ |\pi|\le s,\ \mathrm{check}(F,\pi)$'' is in NP in $(F,1^s)$; $\SAT\in P$ gives polynomial-time search by self-reducibility. Try $s=1,2,4,\dots$ up to the first success (at most $2W_R(F)$ up to the cost/length polynomial), then binary-search the minimum.
\end{proof}
So granting the optimal sequence for free does not change lower bounds by more than a polynomial, in the only world (P$=$NP) where the bridge is being used. It is also literally what the owner's clarification says for the lower-bound direction.

\begin{proposition}[Weak $R$ makes the bridge circular]\status{PROVED/EXACT}
If $\overline W_R$ is unconditionally superpolynomial (e.g.\ $R$ = Resolution, Haken 1985; cutting planes, bounded-depth Frege), then B3$(R)\iff \SAT\notin P$.
\end{proposition}
\begin{proof} The consequent of B3$(R)$ is false, so B3$(R)$ is equivalent to the falsity of its hypothesis. \end{proof}
Hence the bridge is informative only for an $R$ with no known superpolynomial lower bound, and a lower bound for \emph{weaker} subsystems (resolution, XOR/Tseitin-compatible subsystems) carries no weight toward P$\ne$NP. Lower bounds transfer only upward in strength: if $R$ p-simulates $Q$ then $W_R\le \mathrm{poly}(W_Q)$, so a bound for $R$ gives one for $Q$.

\begin{proposition}[Quantifier order]\status{PROVED/EXACT}
(a) $\forall A\,\exists R$: if $A$ is a correct polynomial SAT decider, $R_A$ (``proof $=$ empty, accept iff $A(F)$ rejects'') is sound, polynomially verified, and $\overline W_{R_A}\le \mathrm{poly}$. Hence $\SAT\in P\Rightarrow$ some polybounded system exists (this is the easy half of Cook--Reckhow).\\
(b) B3$(R)$ for a \emph{preselected} $R$ is the $\exists R\,\forall A$ statement; it is equivalent to ``under $\SAT\in P$, $R$ weakly simulates (length-only) every $R_A$, $A$ correct''.
\end{proposition}
(b) is a restatement; its content is that $R$ cannot see which $A$ will be correct, and correctness of $A$ is a $\Pi_2$ fact. B1 of report 377 is the pointwise form of the same statement.

\begin{theorem}[Reflective system: what a lower bound would and would not give]\status{PROVED/EXACT}
Let $T$ be a sound r.e.\ theory. Define $R_T$: a proof of $F$ is $(e,k,\sigma,\tau)$ where $\sigma$ is a $T$-proof of $\mathrm{Corr}(e,k)$ (``machine $e$ decides $\SAT$ in time $\le n^k+k$'') and $\tau$ is the (polynomial-size) run of $M_e$ on $F$ rejecting. $R_T$ is a sound Cook--Reckhow system. If $T\vdash\mathrm{Corr}(e,k)$ for some $(e,k)$ then $\overline W_{R_T}(N)\le c_{T,e}+\mathrm{poly}_k(N)$. Contrapositively, a superpolynomial lower bound for $R_T$ implies $T$ proves no $\mathrm{Corr}(e,k)$ --- an \emph{independence} statement, not $\SAT\notin P$.
\end{theorem}
\begin{proof} Soundness: $\mathrm{Corr}$ true by soundness of $T$, so rejection implies unsat. Size: $|\sigma|$ is a constant for fixed $(e,k)$. \end{proof}
Thus B3$(R_T)$ holds iff $\SAT\in P\Rightarrow\exists(e,k)\ T\vdash \mathrm{Corr}(e,k)$ (``provable P$=$NP''), and the standard Levin-style construction does not supply this: it gives correctness on satisfiable inputs (outputs are verified) but the rejection time bound $k$ must be known and proved. \status{OPEN} for every $T$ considered (not asserted false).

\section{Consequences for the five roles}
\begin{enumerate}
\item \textbf{Lower-bound lead:} a superpolynomial $\overline W_R$ bound for $R$ with no known bound is the \emph{necessary} half; it is not P$\ne$NP-bearing without B3$(R)$. Target $R$ at least as strong as Extended Frege or \texttt{Rmix341}-with-DEF; resolution-type bounds are countercontrols only.
\item \textbf{Bridge builder:} B3$(R)$ $\equiv$ ``P$=$NP $\Rightarrow R$ polybounded''. Any proof must use only P$=$NP, i.e.\ must show a \emph{short $R$-proof of $\neg$SAT for an arbitrary $F$ from the existence of $A$}. Since $R$ cannot be uniform in $A$, progress requires $R$ to internalize reflection (Theorem above) --- the identified crux.
\item \textbf{Hypergraph/certificates:} hard families must resist $R$'s shortcuts (Buss PHP in polynomial Frege, Tseitin via XOR); unconditional bounds against strong $R$ are the open core (Frege/EF lower bounds are open).
\item \textbf{Resource auditor:} Proposition 1 justifies not charging strategy search for lower bounds; charge verification and peak memory of the checker separately.
\item \textbf{Red team:} falsifiers to attack: any claimed B3 proof that never uses an assumption beyond SAT$\in$P must exhibit why it does not also prove B3 for Resolution (Proposition 2 forbids that, since it would prove P$\ne$NP).
\end{enumerate}

\section{Literature check (this run)}
Web search (standard mode) returned background consistent with the above: polynomially bounded systems exist iff NP$=$coNP (Cook--Reckhow); whether an optimal proof system exists is open; P$=$NP makes every proof system automatable (cf.\ Prop.\ 1), automatability is distinct from polynomial boundedness; Atserias--M\"uller, ``The Proof Analysis Problem'' (FOCS 2025; \url{https://arxiv.org/abs/2506.16956}, reported by the search tool as showing EF proof analysis NP-complete; not independently read), and Krajicek--Pudlak (1998) non-automatability of EF under cryptographic assumptions. \status{CANDIDATE} clue: PAP-style hardness on EF shows that \emph{recognizing} what an EF proof achieves is hard, a possible source of a lower-bound handle; no transfer theorem is claimed. No source supported ``EF is polybounded if P$=$NP''; that is the open B3(EF).

\section{Next target}
Decide, for a concrete preselected $R$ (suggested: EF-with-signed-DEF), whether B3 can follow from a formalization of ``$\SAT\in P$'' inside $R$'s own reasoning (bounded-arithmetic route: PV/$S_2^1$ proves correctness only for provably-correct $A$; Cook--Krajicek, Krajicek--Pudlak). First concrete subtask: state precisely which fragment of $\mathrm{Corr}(e,k)$ for an \emph{explicit} Levin-type machine is provable in PV.

\medskip\noindent P versus NP remains OPEN.

% END_OF_VERIFIED_REPORT::CLAUDE::WORSTCASE-BRIDGE-0001
\end{document}
